\subsection{多条记忆：为什么会互相干扰？}
\phantomsection\label{intro:05-memory-interference}
\phantomsection\label{concept:memory-interference}
\phantomsection\label{partq:memory}
对 $T_k$ 条行向量 Key--Value 配对，最直接的写入是外积累加：
\begin{equation}
\boxed{\mathcal{M}=\sum_{i=1}^{T_k}k_i^\top v_i},
\qquad
\boxed{x\mathcal{M}=\sum_{i=1}^{T_k}(xk_i^\top)v_i}.
\end{equation}
这里 $\mathcal{M}:d_k\times d_v$，$x:1\times d_k$，读取输出仍是一行。

\subsubsection{两条记忆的二维例子}
\[
k_A=(1,0),\quad v_A=(1,0),\qquad
k_B=(0.6,0.8),\quad v_B=(0,1).
\]
写入矩阵分别是
\[
k_A^\top v_A=\begin{bmatrix}1&0\\0&0\end{bmatrix},
\qquad
k_B^\top v_B=\begin{bmatrix}0&0.6\\0&0.8\end{bmatrix}.
\]
因此
\[
\boxed{\mathcal{M}_{\rm simple}=\begin{bmatrix}1&0.6\\0&0.8\end{bmatrix}}.
\]
查询 A：
\[
k_A \mathcal{M}_{\rm simple}
=(1,0)\begin{bmatrix}1&0.6\\0&0.8\end{bmatrix}
=\boxed{(1,0.6)}
=\underbrace{(1,0)}_{\text{正确信号}}
+\underbrace{(0,0.6)}_{\text{来自 B 的串扰}}.
\]
B 被激活，是因为 $k_Ak_B^\top=0.6$ 不为零；
两条 Key 虽然都是单位长度，却不垂直。

对单位行向量 Key $k_j$（$k_jk_j^\top=1$），恒有
\[
k_j\mathcal{M}=v_j+\sum_{i\ne j}(k_jk_i^\top)v_i.
\]
其他 Key 的匹配分数不为零时，就会产生干扰。
